%Chapter 4. Monte Carlo design: all definitions and constants match the final six-file
% four-method Python workflow.
\chapter{Monte Carlo Design}
\label{chap:simulation-design}

This chapter specifies the experiment used to compare the four estimators
developed in Chapter~\ref{chap:methods}.  Its purpose is to separate the
scientific design from the numerical findings.  The data-generating process,
component-selection rules, diagnostic thresholds, loss functions, and random
number streams recorded here define the frozen final experiment; none is
redefined using the outcomes reported in the next chapter.  The design retains
the three functional linear models and principal numerical settings used in
the replication study of \textcite{babii_carrasco_tsafack_2026}, while
restricting the comparison to the four estimators relevant to this thesis:
released-code CG--FPLS, raw FPLS, Arnoldi FPLS, and FPCR.

The three models form controlled comparisons.  Model~2 changes only the slope
relative to Model~1, whereas Model~3 changes only the score-variance sequence
that determines the predictor covariance.
Within every Monte Carlo replication, all four estimators receive exactly the
same training and test observations.  Differences between methods within a
model therefore reflect estimation, tuning, and numerical implementation
rather than different simulated samples.

\section{Discretised functional data-generating process}
\label{sec:simulation-dgp}

\subsection{Grid, basis, and inner product}

The functional domain is $\mathcal I=[0,1]$.  Each curve is observed without
additional measurement error on the endpoint-inclusive grid
\begin{equation}
    s_t=\frac{t-1}{T-1},
    \qquad t=1,\ldots,T,
    \qquad T=200.
    \label{eq:simulation-grid}
\end{equation}
All integrals used by the estimators and evaluation criteria are represented
by the equal-weight discrete inner product
\begin{equation}
    \langle f,g\rangle_T
    =\frac1T\sum_{t=1}^{T}f(s_t)g(s_t).
    \label{eq:simulation-inner-product}
\end{equation}
This is the endpoint average used in the code, rather than a trapezoidal rule.
Consequently, every displayed finite-sample formula is interpreted under the
same numerical geometry as the estimators.

The fixed dictionary contains $J=100$ functions.  Following the indexing in
the released replication code, it is
\begin{equation}
    \phi_1(s)=1,
    \qquad
    \phi_j(s)=\sqrt{2}\cos(j\pi s),
    \quad j=2,\ldots,J.
    \label{eq:simulation-basis}
\end{equation}
The indexing in \eqref{eq:simulation-basis} is intentional.  The code first
constructs cosine frequencies $1,\ldots,J$ and then replaces its first column
by the constant function; hence $\sqrt{2}\cos(\pi s)$ is not included.  Stating
this convention explicitly prevents an apparently harmless re-indexing from
creating a different experiment.  In continuous $L^2[0,1]$ the retained
functions are mutually orthonormal.  On the finite endpoint grid their
inner products are evaluated exactly by \eqref{eq:simulation-inner-product}
and are not assumed to equal their continuum values.

\subsection{Predictors and responses}

For observation $i$, draw independent standard normal scores
\begin{equation}
    Z_{ij}\stackrel{\mathrm{iid}}{\sim}\mathcal N(0,1),
    \qquad j=1,\ldots,J,
\end{equation}
and construct the functional predictor
\begin{equation}
    X_i(s_t)
    =\sum_{j=1}^{J}\sqrt{\lambda_j}\,Z_{ij}\phi_j(s_t).
    \label{eq:simulation-predictor}
\end{equation}
Here $\lambda_j$ is exactly the variance multiplier of the independent score
$Z_{ij}$.  In the underlying continuous orthonormal system these quantities
are covariance eigenvalues.  Because the endpoint-grid inner product does not
make the retained dictionary exactly orthonormal, however, they should be read
as score variances---not as the exact eigenvalues of the finite $T\times T$
discrete covariance matrix.  This distinction does not affect any formula
below, which is evaluated directly under the discrete geometry.
The scalar response is generated from
\begin{equation}
    Y_i
    =\langle X_i,\beta\rangle_T+\varepsilon_i
    =\frac1T\sum_{t=1}^{T}X_i(s_t)\beta(s_t)+\varepsilon_i,
    \qquad
    \varepsilon_i\stackrel{\mathrm{iid}}{\sim}\mathcal N(0,\sigma^2),
    \label{eq:simulation-response}
\end{equation}
where $\sigma=1$ and the errors are independent of all scores.  The population
means of $X_i$ and $Y_i$ are therefore zero and the intercept is known to be
zero.  The simulation passes the generated observations directly to all four
estimators; it does not subtract realised sample means.  This common convention
isolates regularisation and numerical effects from an additional centring
step.

For a coefficient vector $b=(b_1,\ldots,b_J)^\top$, the slope evaluated on the
grid is
\begin{equation}
    \beta(s_t)=\sum_{j=1}^{J}b_j\phi_j(s_t).
    \label{eq:simulation-slope-expansion}
\end{equation}
Equations \eqref{eq:simulation-predictor}--\eqref{eq:simulation-slope-expansion}
define the data directly at the $T$ grid points.  There is thus no preliminary
smoothing, basis estimation, or interpolation stage in this controlled study.

\section{The three controlled models}
\label{sec:simulation-models}

Define baseline slope coefficients and score variances by
\begin{equation}
    b_j^{(0)}=\frac{4}{j^{2.7}},
    \qquad
    \lambda_j^{(0)}=\frac{2}{j^{1.1}},
    \qquad j=1,\ldots,J.
    \label{eq:simulation-baseline-sequences}
\end{equation}
Two targeted modifications are then introduced:
\begin{align}
    b_j^{(1)}
    &=
    \begin{cases}
        4, & j\leq5,\\
        4j^{-2.7}, & j>5,
    \end{cases}
    \label{eq:simulation-modified-slope}\\
    \lambda_j^{(1)}
    &=
    \begin{cases}
        2, & j\leq5,\\
        2j^{-1.1}, & j>5.
    \end{cases}
    \label{eq:simulation-modified-spectrum}
\end{align}
The resulting combinations are summarised in
Table~\ref{tab:simulation-model-definitions}.

\begin{table}[htbp]
    \centering
    \small
    \caption[Definitions and roles of the simulation models]{Definition and controlled role of the three simulation models.}
    \label{tab:simulation-model-definitions}
    \begin{tabular}{@{}
        >{\raggedright\arraybackslash}p{1.7cm}
        >{\centering\arraybackslash}p{2.3cm}
        >{\centering\arraybackslash}p{2.3cm}
        >{\raggedright\arraybackslash}p{6.1cm}@{}}
        \toprule
        Model & Slope coefficients & Score variances & Controlled role \\
        \midrule
        Model 1
        & $b^{(0)}$
        & $\lambda^{(0)}$
        & Baseline: rapidly decaying slope coefficients and a power-law
          score-variance sequence. \\

        Model 2
        & $b^{(1)}$
        & $\lambda^{(0)}$
        & Slope-only change: the first five coefficients are set to $4$, so
          coefficients $2$--$5$ increase while the first remains unchanged;
          the predictor distribution is unchanged. \\

        Model 3
        & $b^{(0)}$
        & $\lambda^{(1)}$
        & Predictor-only change: the first five score variances are set to
          $2$, so variances $2$--$5$ increase while the first remains
          unchanged; the true slope is unchanged. \\
        \bottomrule
    \end{tabular}
\end{table}

Figure~\ref{fig:simulation-slope-structure} shows both the true slopes and the
first twenty coefficient magnitudes.  Models~1 and 3 have the same target
function.  Model~2 distributes substantial slope energy across the first five
directions and is therefore more oscillatory.  This distinction tests whether
a supervised method can exploit response-relevant directions differently from
a method ordered only by predictor variance.

\begin{figure}[htbp]
    \centering
    \includegraphics[width=0.99\textwidth]{true_slope_structure.pdf}
    \caption[True slopes and coefficient magnitudes]{True slope functions and coefficient magnitudes in the three
    Monte Carlo models.  The curves use $J=100$ functions and are evaluated on
    the $T=200$ grid.  Only the first twenty coefficients are shown in the
    right-hand panels; all one hundred enter the data-generating process.
    Rows correspond to Models~1--3, with the slope functions on the left.}
    \label{fig:simulation-slope-structure}
\end{figure}

Figure~\ref{fig:simulation-covariance-spectrum} isolates the predictor-side
change.  Models~1 and 2 share the power-law score-variance sequence
$2j^{-1.1}$.  Model~3 has a flat leading block of five score variances followed
by the same tail.  The second panel reports cumulative predictor variation,
not predictive signal: these two quantities need not rank directions in the
same way.

\begin{figure}[htbp]
    \centering
    \includegraphics[width=0.96\textwidth]{covariance_spectrum.pdf}
    \caption[Designed covariance spectra]{Designed covariance structure.  The left panel displays the
    first thirty score variances---the covariance eigenvalues in the
    continuous design---on a logarithmic scale; the right panel shows their
    cumulative share over all $J=100$ components.  The dotted vertical line
    separates the leading block from the common tail.}
    \label{fig:simulation-covariance-spectrum}
\end{figure}

\section{Signal strength and effective directionality}
\label{sec:simulation-signal}

The scale of a slope alone does not determine prediction difficulty.  Let
\begin{equation}
    c_j=\langle\beta,\phi_j\rangle_T,
    \qquad
    q_j=\lambda_jc_j^2.
    \label{eq:simulation-signal-contribution}
\end{equation}
Because the scores are independent and standard normal, $q_j$ is the exact
discrete contribution of component $j$ to the variance of the noiseless
response.  In particular,
\begin{align}
    \mathbb E\|X\|_T^2
    &=\sum_{j=1}^{J}\lambda_j\|\phi_j\|_T^2,
    \label{eq:simulation-predictor-variance}\\
    \operatorname{Var}\{\langle X,\beta\rangle_T\}
    &=\sum_{j=1}^{J}q_j,
    \label{eq:simulation-signal-variance}\\
    \operatorname{Var}(Y)
    &=\sum_{j=1}^{J}q_j+\sigma^2.
    \label{eq:simulation-response-variance}
\end{align}
Appendix Proposition~\ref{prop:app-simulation-variance} proves these identities.
The signal-to-noise ratio is defined as
\begin{equation}
    \operatorname{SNR}
    =\frac{\sum_{j=1}^{J}q_j}{\sigma^2}.
    \label{eq:simulation-snr}
\end{equation}

Table~\ref{tab:simulation-derived-properties} reports values computed from the
finite $J=100$, $T=200$ design itself; they are not Monte Carlo estimates.
Since $\sigma^2=1$, signal variance and SNR are numerically identical.  The
substantially larger signal variance in Model~2 also means that raw prediction
errors should primarily be compared across methods within a model, not used as
an unnormalised measure of which model is intrinsically easiest.

\begin{table}[htbp]
    \centering
    \scriptsize
    \caption[Exact properties of the discrete generating design]{Exact derived properties under the discrete data-generating
    process.  The final column gives the first component counts reaching
    $90\%$, $95\%$, and $99\%$ of predictive signal.}
    \label{tab:simulation-derived-properties}
    \begin{tabular}{@{}lrrrrrc@{}}
        \toprule
        Model
        & $\|\beta\|_T^2$
        & $\mathbb E\|X\|_T^2$
        & $\operatorname{Var}\{\langle X,\beta\rangle_T\}$
        & SNR
        & First-five share
        & $m_{90}/m_{95}/m_{99}$ \\
        \midrule
        Model 1 & 16.484 & 8.589  & 32.520 & 32.520 & 99.986\% & $1/1/2$ \\
        Model 2 & 81.439 & 8.589  & 71.306 & 71.306 & 99.937\% & $4/5/5$ \\
        Model 3 & 16.484 & 14.311 & 33.061 & 33.061 & 99.986\% & $1/1/2$ \\
        \bottomrule
    \end{tabular}
\end{table}

The allocation is made explicit in
Figure~\ref{fig:simulation-signal-allocation}.  Model~2 needs five components
to reach $99\%$ of the exact predictive signal, whereas Models~1 and 3 reach
that level after two.  This statement describes the population design; it does
not predetermine how many empirical components a noisy, data-driven estimator
will select.

\begin{figure}[htbp]
    \centering
    \includegraphics[width=0.96\textwidth]{predictive_signal_allocation.pdf}
    \caption[Predictive-signal allocation]{Componentwise and cumulative predictive-signal allocation,
    calculated as $q_j=\lambda_j\langle\beta,\phi_j\rangle_T^2$ under the
    exact discrete design.  Open markers denote the first component count
    exceeding $99\%$ of total signal.}
    \label{fig:simulation-signal-allocation}
\end{figure}

Representative predictor curves are shown in
Figure~\ref{fig:simulation-predictor-examples}.  Models~1 and 2 have the same
predictor distribution and therefore share a panel.  For this setup
illustration only, the same Gaussian score draws are used in both panels so
that the visibly larger variation for Model~3 can be attributed to its
modified spectrum.  The shaded regions are exact pointwise $90\%$ Gaussian
bands under the corresponding finite data-generating process.  They are not
estimated confidence bands and are not used by any estimator.

\begin{figure}[htbp]
    \centering
    \includegraphics[width=0.96\textwidth]{functional_predictor_examples.pdf}
    \caption[Functional predictors and pointwise population bands]{Representative functional predictors and their exact pointwise
    $90\%$ bands: Models~1 and 2 above, Model~3 below.  The shared score draws make the spectrum change visually
    comparable; the actual Monte Carlo models use independent random-number
    streams.}
    \label{fig:simulation-predictor-examples}
\end{figure}

These population contrasts define the signal each estimator must recover.
The Monte Carlo protocol now determines how sampling variation and
data-driven tuning act on those contrasts.

\section{Monte Carlo protocol}
\label{sec:simulation-protocol}

The full experiment uses $R=5{,}000$ independent replications for each of the
three models.  Within replication $r$ of a given model, the following protocol
is applied.
\begin{enumerate}
    \item Generate a training sample of $n=100$ curve--response pairs from
    \eqref{eq:simulation-predictor}--\eqref{eq:simulation-response}.

    \item Construct the empirical moments
    \begin{equation}
        \widehat K=\frac{X^\top X}{nT},
        \qquad
        \widehat r=\frac{X^\top y}{n},
        \label{eq:simulation-empirical-moments}
    \end{equation}
    and fit all four methods to the same training sample.

    \item Generate an independent test sample of $n_{\mathrm{test}}=100$ from
    the same model.  No test observation is used for tuning.

    \item Store the selected dimension, slope and prediction errors, and
    method-specific numerical diagnostics.
\end{enumerate}

The candidate set is $m\in\{1,\ldots,70\}$ for every method.  Thus $m=0$ is
never reported as an estimator.  For CG--FPLS, absence of a threshold crossing
returns the $m=70$ estimate.  The common upper bound keeps the candidate
complexities comparable while remaining below the empirical rank limit; the
point selected along that path is nevertheless method-specific.  Methods are
paired within a replication, but the three model streams are independent.

\section{Estimator-specific tuning}
\label{sec:simulation-tuning}

The method labels refer to complete procedures, including their own tuning
rules.  Table~\ref{tab:simulation-method-tuning} provides an overview; the
details below fix the criteria and constants.

\begin{table}[htbp]
    \centering
    \small
    \caption[Estimators and tuning mechanisms]{The four reported estimators and their tuning mechanisms.}
    \label{tab:simulation-method-tuning}
    \begin{tabular}{@{}
        >{\raggedright\arraybackslash}p{2.1cm}
        >{\raggedright\arraybackslash}p{2.6cm}
        >{\raggedright\arraybackslash}p{3.2cm}
        >{\raggedright\arraybackslash}p{3.8cm}@{}}
        \toprule
        Method & Restricted problem & Selection of $m$ & Numerical construction \\
        \midrule
        CG--FPLS
        & Moment-residual least squares
        & First discrepancy-threshold crossing; $\tau=1.01$, $\delta=0.1$
        & Released short recurrence; internal FPCR noise pilot \\

        Raw FPLS
        & Response least squares on $\widehat{\mathcal K}_m$
        & Five-fold prediction cross-validation
        & Unscaled covariance powers and normal equations \\

        Arnoldi FPLS
        & Same response least-squares problem as raw FPLS
        & Five-fold prediction cross-validation
        & CGS2 Arnoldi basis and response-space QR; rank tolerance $10^{-10}$ \\

        FPCR
        & Response least squares on leading empirical eigenvectors
        & Response-space GCV
        & Symmetric eigendecomposition and spectral cut-off \\
        \bottomrule
    \end{tabular}
\end{table}

\subsection{Five-fold tuning of raw and Arnoldi FPLS}

One random partition $I_1,\ldots,I_5$ is generated in each replication and is
used by both FPLS implementations.  For method $a\in\{\mathrm{Raw},
\mathrm{Arnoldi}\}$, let $\widehat\beta_{m,a}^{(-k)}$ be the estimate fitted
without fold $I_k$.  Its criterion is
\begin{equation}
    \operatorname{CV}_{a}(m)
    =\frac1n\sum_{k=1}^{5}\sum_{i\in I_k}
      \left\{Y_i-\langle X_i,
      \widehat\beta_{m,a}^{(-k)}\rangle_T\right\}^{2}.
    \label{eq:simulation-fpls-cv}
\end{equation}
Each method minimises its own curve.  The common folds remove avoidable Monte
Carlo variation, but separate minimisation allows finite-precision instability
in the raw basis to affect the complete procedure honestly.  A raw component
count is eligible only when its fit is finite and valid in the full sample and
in every training fold.  Arnoldi uses two-pass classical Gram--Schmidt (CGS2)
with rank tolerance $10^{-10}$.  If several eligible values attain the same
minimum numerically, the smallest is retained.

\subsection{Response-GCV tuning of FPCR}

FPCR uses the response-space criterion already introduced in
\eqref{eq:methods-response-gcv}:
\begin{equation}
    \operatorname{GCV}_{Y}(m)
    =\frac{n^{-1}\|y-X\widehat\beta_m^{\mathrm{FPCR}}/T\|_2^2}
           {(1-m/n)^2}.
    \label{eq:simulation-fpcr-gcv}
\end{equation}
Only numerically positive empirical eigenvalues are eligible.  Under the
present $m_{\max}=70<n-1$ design, the shared upper bound is the active path
limit.  The smallest minimiser of \eqref{eq:simulation-fpcr-gcv} is selected.

\subsection{Discrepancy tuning of CG--FPLS}

For CG--FPLS, define the discrete moment norm and predictor-energy estimate by
\begin{equation}
    \|e_m\|_T
    =\left(\frac1T e_m^\top e_m\right)^{1/2},
    \qquad
    \widehat E_X
    =\frac1n\sum_{i=1}^{n}\frac1T X_i^\top X_i,
    \label{eq:simulation-cg-norms}
\end{equation}
where $e_m=\widehat r-\widehat K\widehat\beta_m^{\mathrm{CG}}$.  The noise
pilot is constructed from the same FPCR path but is selected internally by
moment-space GCV,
\begin{equation}
    \operatorname{GCV}_{M}(q)
    =\frac{T^{-1}\|\widehat r-\widehat K
      \widehat\beta_q^{\mathrm{FPCR}}\|_2^2}
           {(1-q/T)^2}.
    \label{eq:simulation-moment-gcv}
\end{equation}
If $\widehat q$ minimises this criterion, then
\begin{equation}
    \widehat\sigma^2
    =\frac1n\|y-X\widehat\beta_{\widehat q}^{\mathrm{FPCR}}/T\|_2^2.
    \label{eq:simulation-cg-sigma-pilot}
\end{equation}
This moment-GCV calculation is used only to estimate the threshold scale; it
is neither the reported FPCR fit nor a fifth comparator.

For interpretation, an ideal version of the discrepancy scale would replace
the unknown quantities below by the population error variance and predictor
energy.  The implemented rule is feasible instead: it uses the internally
estimated $\widehat\sigma^2$ in
\eqref{eq:simulation-cg-sigma-pilot} and the sample energy
$\widehat E_X$.  Thus $\widehat\nu_n$ is a plug-in threshold, not an oracle
threshold supplied with the true $\sigma^2$ or the population covariance.
The selected component count is the first $m\geq1$ satisfying
\begin{equation}
    \|e_m\|_T
    \leq
    \widehat\nu_n
    =1.01\left(
      \frac{2\widehat\sigma^2\widehat E_X}{0.1n}
      \right)^{1/2}.
    \label{eq:simulation-cg-threshold}
\end{equation}
If no iterate crosses the threshold, $m=70$ is used and the absence of a
crossing is recorded.  Together with the cross-validation and GCV rules above,
this defines the four complete procedures whose consequences are evaluated
through slope and prediction loss.  The rule in
\eqref{eq:simulation-cg-threshold} is the single CG selection procedure in the
final workflow; no oracle or supplementary CG variant is included.

\section{Performance measures and summaries}
\label{sec:simulation-criteria}

\subsection{Slope estimation and independent prediction}

For method $a$, the integrated squared error in one replication is the
discrete loss
\begin{equation}
    \operatorname{ISE}_{a}
    =\|\widehat\beta_a-\beta\|_T^2
    =\frac1T\sum_{t=1}^{T}
      \{\widehat\beta_a(s_t)-\beta(s_t)\}^2.
    \label{eq:simulation-ise}
\end{equation}
The independent test mean squared prediction error is
\begin{equation}
    \operatorname{MSPE}_{a}
    =\frac1{n_{\mathrm{test}}}
      \sum_{i=1}^{n_{\mathrm{test}}}
      \{Y_i^{\mathrm{test}}
      -\langle X_i^{\mathrm{test}},\widehat\beta_a\rangle_T\}^2.
    \label{eq:simulation-mspe}
\end{equation}
MSPE includes irreducible response noise.  Conditional on a training-sample
estimate and averaging over a new test observation,
\begin{equation}
    \mathbb E\{\operatorname{MSPE}_{a}\mid\widehat\beta_a\}
    =\sigma^2+
      \sum_{j=1}^{J}\lambda_j
      \langle\phi_j,\widehat\beta_a-\beta\rangle_T^2.
    \label{eq:simulation-mspe-risk}
\end{equation}
Thus the oracle noise floor is $\sigma^2=1$, while the second term is excess
prediction risk.  Equation \eqref{eq:simulation-mspe-risk} is the discrete
counterpart of Appendix Proposition~\ref{prop:app-risk-identity}.

For every model--method pair, the archive reports the mean and its Monte Carlo
standard error, the median, the $90$th and $99$th percentiles, and the maximum
of both ISE and MSPE.  If $L_1,\ldots,L_{R_f}$ are the finite values of a loss,
the reported standard error is
\begin{equation}
    \operatorname{MCSE}(\overline L)
    =\frac{s_L}{\sqrt{R_f}},
    \label{eq:simulation-mcse}
\end{equation}
where $s_L$ is their sample standard deviation.  Medians and upper quantiles
are reported alongside means because rare numerical failures can make the
error distributions extremely right-skewed.

\subsection{Bias--variance decomposition}

The deterministic replay stage reconstructs each selected slope whenever the
fit is numerically reproducible.  Unstable raw fits are omitted only from
replay-based function-valued diagnostics---the bias--variance decomposition,
basis-coefficient distributions, pointwise distributions, and normal
quantile--quantile plots.  They remain in every primary ISE and MSPE summary.
With
$\overline\beta_a=R^{-1}\sum_{r=1}^{R}\widehat\beta_{a,r}$, the empirical
slope decomposition is
\begin{equation}
    \frac1R\sum_{r=1}^{R}
      \|\widehat\beta_{a,r}-\beta\|_T^2
    =\|\overline\beta_a-\beta\|_T^2
     +\frac1R\sum_{r=1}^{R}
       \|\widehat\beta_{a,r}-\overline\beta_a\|_T^2.
    \label{eq:simulation-bias-variance}
\end{equation}
The two right-hand terms are labelled squared bias and variance.  Prediction
error is decomposed analogously after weighting the coordinatewise bias and
variance by $\lambda_j$, with $\sigma^2$ added as irreducible noise.
Appendix Proposition~\ref{prop:app-mc-decomposition} proves the finite-$R$
identity.  The factor $R^{-1}$, rather than $(R-1)^{-1}$, is used so the two
terms add exactly to the Monte Carlo mean squared error.

The primary risk summaries retain every finite replication.  A secondary
paired raw-stable subset is used only for a like-for-like bias--variance
display when the selected raw power-basis fit can be reconstructed reliably.
Whenever that subset is used, the same replications are retained for all four
methods and the restriction is labelled explicitly.  It does not replace the
full-sample ISE or MSPE comparison.

\subsection{Complexity, instability, and tail events}

Selected component count is recorded for every fit and summarised by its mean
and median.  Complexity is not itself an accuracy measure, but it shows how
aggressively each tuning rule regularises the inverse problem.

The simulation also distinguishes statistical error from numerical failure.
For raw FPLS, let $G_m$ be the raw Krylov power matrix and
$Z_m=XG_m/T$ its response-space design.  Let $\widetilde G_m$ and
$\widetilde Z_m$ denote the same matrices after normalising every column to
unit Euclidean length.  The diagnostic is
\begin{equation}
    \Delta_m^{\mathrm{Raw}}
    =\max\{\kappa_2(\widetilde G_m),\kappa_2(\widetilde Z_m)\}.
    \label{eq:simulation-raw-condition}
\end{equation}
The column normalisation removes trivial scale growth and focuses the measure
on near-collinearity.  A selected raw fit is flagged as numerically unstable
when its loss is non-finite or
\begin{equation}
    \Delta_m^{\mathrm{Raw}}
    \geq\varepsilon_{\mathrm{mach}}^{-1/2}
    \approx 6.71\times10^7.
    \label{eq:simulation-raw-threshold}
\end{equation}

For Arnoldi FPLS the corresponding diagnostic is the spectral-norm
orthogonality defect
\begin{equation}
    \Delta_m^{\mathrm{Arnoldi}}
    =\|Q_m^\top Q_m-I_m\|_2.
    \label{eq:simulation-arnoldi-defect}
\end{equation}
A fit is flagged when its loss is non-finite or the defect is at least
$\varepsilon_{\mathrm{mach}}^{1/2}\approx1.49\times10^{-8}$.  CG--FPLS and
FPCR are flagged only if ISE or MSPE is non-finite because their relevant path
diagnostics are already governed by the recurrence and positive-eigenvalue
checks.

Finally, a finite ISE or MSPE value is marked as an extreme tail event if it
exceeds $100$ times the positive median for its own model--method pair.  This
threshold is descriptive rather than a deletion rule: extreme observations
remain in means, maxima, archives, and full-distribution figures.  Reporting
failure rates, instability rates, and tail counts alongside robust summaries
prevents a numerically fragile procedure from appearing reliable merely
because its typical replication is accurate.  Because these classifications
must themselves be reproducible, the random streams and replay rules are fixed
as carefully as the losses and thresholds.

\section{Randomness and reproducibility}
\label{sec:simulation-reproducibility}

The top-level seed is $2026$.  A NumPy \texttt{SeedSequence} deterministically
spawns one independent child stream for each model.  Within a model, the
training sample is generated first, followed by the random five-fold
partition and then the independent test sample.  The order matters for exact
replay.  The four methods share the training and test data, and raw and
Arnoldi FPLS additionally share the folds.  Models do not share their realised
samples.

The final plotting stage deterministically regenerates selected estimates
from the seed, model definition, and archived component counts before checking
their ISE and MSPE against the frozen archive.  This replay validation is an
audit of internal consistency, not a second simulation or an exclusion rule.
The exact command, file manifest, validation tolerances, and environment fields
to record are kept in the accompanying reproducibility documentation rather
than in the statistical argument.

Exact reproduction of the stabilised methods is expected on a common software
stack.  Small platform differences are more plausible for raw FPLS because
ill-conditioned linear algebra is deliberately part of the method being
diagnosed; the archived configuration and validation report make those
differences visible.

Because the models isolate slope and covariance changes while every method
receives the same samples, the results can separate statistical regularisation
from numerical implementation.  The fixed tuning and replay rules also ensure
that adverse tail behaviour remains part of the comparison.
Chapter~\ref{chap:estimation-results} examines how those mechanisms appear in
slope recovery and prediction.
